\documentclass[12pt]{article}
\usepackage{hyperref, amsmath, amssymb, xcolor}
\usepackage[margin=1in]{geometry}
\usepackage{natbib}
\newcommand{\mnras}{MNRAS}
\newcommand{\apj}{ApJ}
\newcommand{\aap}{A\&A}
\newcommand{\apjl}{ApJL}
\newcommand{\aj}{AJ}

\newcommand{\referee}[1]{\medskip\noindent\textit{\textcolor{gray}{Referee: #1}}\medskip}
\newcommand{\response}[1]{\noindent #1\medskip}

\begin{document}

\noindent Dear Referee,

\medskip

We thank you for your careful reading and constructive comments. We have revised the
manuscript to address every point raised. Changes in the manuscript are marked using the
AASTeX \texttt{trackchanges} commands (\verb|\added|, \verb|\deleted|, \verb|\replaced|)
and are highlighted in the compiled PDF.

\bigskip
\hrule
\bigskip

%% -----------------------------------------------------------------------
\referee{The manuscript is well organized and straightforward, although there is
a problem on page 3, lines 86--101, and lines 60--85 appear to be nearly identical
text. This should be consolidated into a single paragraph.}

\response{We agree. Lines 86--101 were an inadvertent duplication of the earlier
paragraph (lines 60--85) describing the comparison between SATLAS and PHOENIX stellar
atmosphere models. The duplicate paragraph has been deleted; only the first, more
complete version is retained.}

\bigskip
%% -----------------------------------------------------------------------
\referee{Lines 151--156, timing synchronization and baseline measurements are generally
not the leading limitation of SII right now. Timing is easily synchronized to better
than 100~psec over baselines up to 10's of km through White Rabbit chipset, and the
telescope separations can be measured to better than 1~cm using inexpensive RTK systems.
The leading limitations in the use of IACT telescopes for SII are: (1)~poor PSF/large
pixels; (2)~fast focal-ratio optics causing interference filter bandpass shifts with
incidence angle; (3)~the finite size of the mirrors; (4)~combining extended measurements
taken over multiple days.}

\response{We thank the referee for this important correction based on hands-on experience
with IACT-based SII systems. We have revised the relevant paragraph to remove clock
synchronization and baseline determination from the list of leading limitations---noting
instead that White Rabbit electronics achieves $<100$~ps inter-telescope synchronization
and RTK-GPS achieves $<1$~cm baseline determination, so neither is a current bottleneck.
We have replaced these with the referee's more accurate characterization of the four
leading IACT SII limitations. We offer the following specific remarks on each:

\medskip
\noindent\textbf{(1) Poor PSF / large pixels.}
We agree that the large IACT camera pixels dilute stellar flux into the sky background,
biasing the $g^{(2)}$ function, and that this is a genuine leading limitation.
We note, however, that the recent H.E.S.S.\ and MAGIC SII campaigns have demonstrated
that external reimaging optics mounted at the IACT focal plane effectively bypass the
native camera, substantially reducing the sky-background contamination without requiring
optical-quality primary mirrors
\citep{2024MNRAS.529.4387A,2024MNRAS.52712243Z,2025MNRAS.537.2334V}.
We have updated the text to acknowledge both the limitation and this practical mitigation.

\medskip
\noindent\textbf{(2) Filter bandpass shift with incidence angle.}
We agree this is a significant concern, especially for the CTAO SSTs which use
Schwarzschild--Couder dual-mirror optics with a very fast effective focal ratio
($f/\sim 0.5$).  For $f/0.5$ the maximum convergent beam half-angle at the focal plane
is $\arctan(D/2f) = \arctan(1) \approx 45^\circ$; photons from the outer annulus of the
primary arrive at the filter at $\sim 45^\circ$ from normal, producing a bandpass blueshift
of $\sim 7\%$ (for a filter effective index $n_\mathrm{eff} \approx 1.9$).
This smears the effective passband across the aperture and can increase the background
contribution from out-of-band photons near the mirror edge, as the referee notes.
We have added text to this effect.
We note that the referee's cited upper bound of $65^\circ$ would require $f/0.23$ optics,
which is beyond the range of practical optical systems; we believe the concern is
well-described by the $\sim 45^\circ$ maximum arising from the nominal $f/0.5$ design.
This does not alter the referee's substantive point, and we have not emphasized this
difference in the paper text.

\medskip
\noindent\textbf{(3) Finite mirror size.}
We agree that the finite mirror diameter distorts the visibility function when the first
visibility null falls within the baseline spanned by the mirror, as for large
($\gtrsim 3$--4~mas) stars with 10~m-class telescopes.
For our specific Red Clump science targets ($\theta \approx 1$--2~mas) and the CTAO SSTs
(4~m primary diameter), the first null of $|\mathcal{V}|^2$ in the $H$-band falls at
baselines of $\sim 150$--300~m, so the visibility variation across a single 4~m mirror is
of order $(4/200)^2 \sim 0.04\%$---negligible for our application.
We have acknowledged the referee's concern in the text and noted this mitigation for our
specific targets.

\medskip
\noindent\textbf{(4) Multi-night background systematics.}
We agree that combining data from many sessions spread over days or months can allow
night-to-night background variations to dominate the noise budget, making the effective
S/N scale worse than $1/\sqrt{T}$.  We have added a sentence noting that this is precisely
why simultaneous multi-baseline arrays such as the CTAO (666 independent baselines from 37
SSTs) are advantageous: the required signal can be accumulated within a single observing
session, avoiding the multi-night systematics problem entirely.}

\bigskip
%% -----------------------------------------------------------------------
\referee{Line 210: ``generated my'' $\to$ ``generated by''.}

\response{Corrected.}

\bigskip
%% -----------------------------------------------------------------------
\referee{Line 234: $T_\mathrm{obs} = 1$~hour. This differs from the figure and table
captions, which state $T_\mathrm{obs} = 2$~hours.}

\response{The figure and table captions give the correct value.
The text stated $T_\mathrm{obs}=1$~h due to an editing error; we have corrected it to
$T_\mathrm{obs}=2$~h to match the captions throughout.}

\bigskip
%% -----------------------------------------------------------------------
\referee{Line 307: having a larger array of telescopes operating simultaneously is
substantially better than having fewer telescopes operate for longer. This is the issue
of controlling systematic shifts in background light over many nights, weeks, or months.}

\response{We agree, and this point reinforces our argument for the CTAO multiplex
advantage.  We have added an explicit sentence at the relevant location noting that
simultaneous multi-baseline observations are preferred over extended single-pair
campaigns specifically because night-to-night background variations can prevent the
S/N from scaling as $1/\sqrt{T}$ in the latter case.}

\bigskip
\hrule
\bigskip

\noindent We believe these revisions fully address the referee's concerns and strengthen
the paper. We appreciate the referee's expert knowledge of practical IACT-SII limitations
and thank them for improving the accuracy of the manuscript.

\bigskip
\noindent Sincerely,

\noindent Alex G.\ Kim and Robin Kaiser

\bibliographystyle{aasjournalv7}
\bibliography{main}

\end{document}
